%%%%%%%%%%%%%%%%%%%%%%%%%%%%%%%%%%%%%%%%%
% Cesar Arguello -- Curriculum Vitae
%
% Requires resume.cls (Trey Hunner, http://www.treyhunner.com/) in the same
% directory, with the custom rSubsection_ra / _ta / _Conference / _Journal /
% _Poster environments already added.
%
% Compile with pdfLaTeX. If you switch to XeLaTeX or LuaLaTeX, delete the
% \usepackage[utf8]{inputenc} line below.
%%%%%%%%%%%%%%%%%%%%%%%%%%%%%%%%%%%%%%%%%

\documentclass{resume}

%----------------------------------------------------------------------------------------
%   FONTS AND ENCODING
%   Loading T1 + Latin Modern avoids Type 3 bitmap font substitution, which is the
%   most common reason a Computer Modern PDF renders heavy/smudged on screen.
%----------------------------------------------------------------------------------------
\usepackage[T1]{fontenc}
\usepackage[utf8]{inputenc}
\usepackage{lmodern}
\usepackage{microtype}

%----------------------------------------------------------------------------------------
%   LAYOUT
%----------------------------------------------------------------------------------------
\usepackage[left=0.7in,top=0.4in,right=0.7in,bottom=0.6in]{geometry}
\usepackage{enumitem}
\usepackage{xcolor}


%----------------------------------------------------------------------------------------
%   LINKS (hyperref last)
%----------------------------------------------------------------------------------------
\usepackage{hyperref}
\definecolor{linknavy}{HTML}{14416E}
\hypersetup{
  colorlinks=true,
  linkcolor=linknavy,
  urlcolor=linknavy,
  citecolor=linknavy,
  pdfauthor={Cesar Arguello},
  pdftitle={Cesar Arguello -- Curriculum Vitae}
}

%----------------------------------------------------------------------------------------
%   HELPERS
%----------------------------------------------------------------------------------------
% resume.cls always opens a bullet list inside rSubsection. Use \nobullets to
% close it cleanly when a block has no bullet points, instead of repeating the
% \vspace{-.65cm}\item[] hack everywhere.
\newcommand{\nobullets}{\item[]\vspace{-1.7em}}


%----------------------------------------------------------------------------------------
%   DOCUMENT START
%----------------------------------------------------------------------------------------
\name{Cesar Arguello}
\address{%
  \href{mailto:cesar.n.arguello.martinez.gr@dartmouth.edu}{cesar.n.arguello.martinez.gr@dartmouth.edu} \\
  \href{https://www.carguellom.com}{carguellom.com} \quad\textbar\quad
  \href{https://github.com/carguelloM}{github.com/carguelloM}%
  % \quad\textbar\quad \href{https://scholar.google.com/citations?user=XXXX}{Google Scholar}
  % \quad\textbar\quad \href{https://orcid.org/0000-0000-0000-0000}{ORCID}
}

\begin{document}

%----------------------------------------------------------------------------------------
%   EDUCATION
%----------------------------------------------------------------------------------------
\begin{rSection}{Education}

\begin{rSubsection}{Ph.D. in Computer Science}{September 2022 -- Present}{Dartmouth College}{Hanover, NH, USA}
  \item[] \textbf{Advisor:} Dr.~David Kotz \\
          \textbf{Research area:} Security and privacy in the lifecycle of IoT for consumer environments
\end{rSubsection}

\begin{rSubsection}{B.S. in Physics \& Computer Science}{August 2018 -- May 2022}{University of Florida}{Gainesville, FL, USA}
  \nobullets
\end{rSubsection}

\begin{rSubsection}{International Baccalaureate Diploma Programme}{August 2016 -- May 2018}{UWC Red Cross Nordic}{Flekke, Norway}
  \nobullets
\end{rSubsection}

\end{rSection}

%----------------------------------------------------------------------------------------
%   PUBLICATIONS
%----------------------------------------------------------------------------------------
\begin{rSection}{Publications}

  \begin{rSubsection_Conference}{Conference Papers}
    \item Ravindra Mangar, \textbf{Cesar Arguello}, David Inyangson, Tina Pavlovich, Karen Gareis, and Tushar Jois. Engaging students from under-represented groups to pursue graduate school in computer science and engineering. \textit{Proceedings of the 56th ACM Technical Symposium on Computer Science Education (SIGCSE TS)}, 2025. Acceptance rate 33\%.
    \item \textbf{Cesar Arguello}, Beatrice Perez, Timothy J. Pierson, and David Kotz. Detecting battery cells with harmonic radar. \textit{Proceedings of the ACM Conference on Security and Privacy in Wireless and Mobile Networks (WiSec)}, 2024. Acceptance rate 21\%.
    \item Beatrice Perez, \textbf{Cesar Arguello}, Timothy J. Pierson, Gregory Mazzaro, and David Kotz. Evaluating the practical range of harmonic radar to detect smart electronics. \textit{Proceedings of the IEEE Military Communications Conference (MILCOM)}, 2023. Acceptance rate 40\%.
  \end{rSubsection_Conference}

  \begin{rSubsection_Journal}{Journal Articles}
    \item Timothy J. Pierson, \textbf{Cesar Arguello}, Beatrice Perez, Wondimu Zegeye, Kevin Kornegay, Carl Gunter, and David Kotz. We need a ``building inspector for IoT'' when smart homes are sold. \textit{IEEE Security \& Privacy}, 2024.
  \end{rSubsection_Journal}

  % Optional: uncomment if you want work-in-progress visible to hiring/fellowship
  % committees. Keep the venue name but do not imply acceptance.
  % \begin{rSubsection_Conference}{Manuscripts Under Review}
  %   \item \textbf{Cesar Arguello}, \ldots. <Paper title>. Under major revision, \textit{Network and Distributed System Security Symposium (NDSS)}, 2027.
  % \end{rSubsection_Conference}

  \begin{rSubsection_Poster}{Posters}
    \item \textbf{Cesar Arguello}, Hunter Searle, Sara Rampazzi, and Kevin Butler. [Poster] A practical methodology for ML-based EM side-channel disassemblers. \textit{Proceedings of the Poster Session of the 7th IEEE European Symposium on Security and Privacy (EuroS\&P)}, 2022.
  \end{rSubsection_Poster}

\end{rSection}

%----------------------------------------------------------------------------------------
%   PATENTS
%----------------------------------------------------------------------------------------
\begin{rSection}{Patents}
  \begin{itemize}[leftmargin=2em, label=$\cdot$, topsep=0pt, itemsep=0.2em, parsep=0pt]
    \item \textbf{Cesar Arguello}, Beatrice Perez, Timothy J. Pierson, and David Kotz. Harmonic radar for battery detection. U.S. Patent Application US19/216,352. Filed May 22, 2025. \textit{(Pending.)}
  \end{itemize}
\end{rSection}

%----------------------------------------------------------------------------------------
%   RESEARCH EXPERIENCE
%----------------------------------------------------------------------------------------
\begin{rSection}{Research Experience}

  \begin{rSubsection_ra}{Graduate Research Assistant}{January 2023 -- Present}{Dartmouth College}{Hanover, NH, USA}{Department of Computer Science}{Dr.~David Kotz and Dr.~Timothy Pierson}
    \item Design, implement, and debug embedded systems across the hardware--software stack in C, C++, Python, and Verilog, building real-time wireless systems on software-defined radio platforms.
    \item Design and implement experiments and circuit prototypes for novel applications of harmonic radar in IoT security, spanning theoretical modeling through field measurement.
    \item Develop research plans targeting conference and journal publication, and review and synthesize the IoT security literature to position new work.
  \end{rSubsection_ra}

  \begin{rSubsection_ra}{Research Assistant}{June 2021 -- May 2022}{University of Florida}{Gainesville, FL, USA}{Florida Institute for Cybersecurity Research}{Dr.~Sara Rampazzi and Dr.~Kevin Butler}
    \item Designed and ran experiments to test theoretical frameworks for EM side-channel disassembly.
    \item Collected, processed, and analyzed EM traces using statistical and machine learning models.
    \item Developed research plans leading to project deliverables, including a peer-reviewed poster.
  \end{rSubsection_ra}

  \begin{rSubsection_ra}{Research Assistant}{May 2021 -- May 2022}{University of Florida}{Gainesville, FL, USA}{Department of Physics, SuperCDMS HVeV-DMC Data Analysis and Simulations Team}{Dr.~Tarek Saab}
    \item Built Ionization Measurement with Phonons At Cryogenic Temperatures (IMPACT) simulation files for the Geant4 simulator.
    \item Analyzed Geant4 simulation output in Python.
    \item Led group discussions on simulation results.
  \end{rSubsection_ra}

  \begin{rSubsection_ra}{Research Assistant}{August 2020 -- February 2021}{University of Florida}{Gainesville, FL, USA}{Department of Mechanical and Aerospace Engineering, ADAMUS Lab}{Dr.~Riccardo Bevilacqua}
    \item Designed, developed, and deployed C++ libraries controlling a CubeSat's magnetometers, EPS, and battery over I\textsuperscript{2}C, integrated with a BeagleBone Black microcontroller.
    \item Tested and debugged flight component prototypes.
    \item Authored weekly reports on flight software progress, blockers, and design decisions.
  \end{rSubsection_ra}

\end{rSection}

%----------------------------------------------------------------------------------------
%   TEACHING EXPERIENCE
%----------------------------------------------------------------------------------------
\begin{rSection}{Teaching Experience}

  \begin{rSubsection_ta}{Graduate Teaching Assistant}{September 2022 -- January 2023}{Dartmouth College}{Hanover, NH, USA}{CS 50: Software Design and Implementation}
    \item Built and maintained automated grading scripts in Bash and Python for student assignments.
    \item Conducted code reviews for programming projects, giving feedback on correctness, efficiency, and style.
    \item Held office hours providing one-on-one mentorship on course material and projects.
  \end{rSubsection_ta}

  \begin{rSubsection_ta}{Undergraduate Teaching Assistant}{January 2021 -- May 2021}{University of Florida}{Gainesville, FL, USA}{CDA 3101: Introduction to Computer Organization}
    \item Led discussion sections on lecture material to help students meet course learning objectives.
    \item Managed assignment evaluation and provided detailed written feedback.
    \item Held office hours to review course material and answer questions on assignments.
  \end{rSubsection_ta}

\end{rSection}

%----------------------------------------------------------------------------------------
%   INDUSTRY EXPERIENCE
%----------------------------------------------------------------------------------------
\begin{rSection}{Industry Experience}

\begin{rSubsection}{Product Development Engineer Intern}{January 2022 -- May 2022}{Intel Corporation}{Santa Clara, CA, USA}
  \item Designed, developed, and debugged sort test programs for server products.
  \item Tested, validated, modified, and redesigned circuits to guarantee component margin to specification.
  \item Evaluated component specification against measured performance to match component requirements with production equipment capability.
\end{rSubsection}

\end{rSection}

%----------------------------------------------------------------------------------------
%   SERVICE
%----------------------------------------------------------------------------------------
\begin{rSection}{Service}
\textbf{Artifact Evaluation Committee}, Network and Distributed System Security Symposium (NDSS) 2027 \\
\textbf{External Reviewer}, IEEE GLOBECOM 2026
\end{rSection}

%----------------------------------------------------------------------------------------
%   HONORS AND AWARDS
%----------------------------------------------------------------------------------------
\begin{rSection}{Honors \& Awards}
\textbf{\textit{Cum Laude}}, B.S., University of Florida \hfill \textit{May 2022} \\
\textbf{Phi Beta Kappa}, Honor Society \hfill \textit{May 2022} \\
\textbf{Davis United World College Scholar} \hfill \textit{August 2018}
\end{rSection}

%----------------------------------------------------------------------------------------
%   TECHNICAL SKILLS
%----------------------------------------------------------------------------------------
\begin{rSection}{Technical Skills}

\begin{tabular}{@{} >{\bfseries}p{1.35in} @{\hspace{2ex}} p{\dimexpr\linewidth-1.35in-2ex\relax} @{}}
Programming       & Python, C, C++, Verilog, MATLAB, Bash \\
Wireless \& SDR   & GNU Radio, USRP, HackRF, ADALM-Pluto/AntSDR, gr-ieee802-11, openwifi, 802.11 PHY/MAC \\
Security tooling  & Wireshark, tcpdump, Scapy, aircrack-ng, hostapd, wpa\_supplicant \\
Libraries         & PyTorch, NumPy, SciPy, scikit-learn, Matplotlib \\
Tools             & Git, GDB, Vivado, Docker, Linux, \LaTeX{} \\
Spoken languages  & English, Spanish \\
\end{tabular}

\end{rSection}

%----------------------------------------------------------------------------------------
%   OPTIONAL SECTIONS (kept commented for later use)
%----------------------------------------------------------------------------------------
% \begin{rSection}{Projects}
% \begin{rSubsection_project}{SDR FM Radio Demodulator}{August 2024}
% Real-time FM radio demodulator for HackRF and RTL-SDR written in Python. \underline{Future work}: support additional SDRs and improve audio streaming. [\href{https://github.com/carguelloM/sdr-fm-radio}{GitHub}]
% \end{rSubsection_project}
%
% \begin{rSubsection_project}{Bi-directional Intercom}{November 2021}
% Bi-directional intercom built from operational amplifiers, BJT transistors, and diodes. [LTspice simulation]
% \end{rSubsection_project}
%
% \begin{rSubsection_project}{Shell for UNIX-Based Systems}{January 2021}
% Shell for Unix-based systems supporting I/O redirection, piping, wildcarding, tilde expansion, and background processing. [\href{https://github.com/carguelloM/CPshell}{GitHub}]
% \end{rSubsection_project}
% \end{rSection}

% \begin{rSection}{Other Academic Activities}
% \textbf{Student Participant}, CyberTruck Challenge \hfill \textit{June 2024}
% \end{rSection}

% \begin{rSection}{Other Experience}
% \begin{rSubsection}{Audio Visual Specialist}{May 2019 -- March 2022}{J. Wayne Reitz Union}{Gainesville, FL, USA}
%   \item Managed and installed audio-visual equipment for shows, meetings, and events.
%   \item Assisted customers with audio-visual technical difficulties.
% \end{rSubsection}
% \end{rSection}

\end{document}